\subsection{The six-type filling construction}\label{app:point-filling}
This proves \cref{prop:point-fill-sharp}. The retained-relator, doubling and rewriting method is credited to~\cite[Appendix~A and Remark~A.3]{DM-H3}; see~\cite[Section~4.2]{Lee12} for the broader filling framework. The choice of six types and the raw block order below give the stated matrix.
\begin{proof}
All areas in this proof are ultimately measured in the original five Johnson relators. No relator is added to the game.

Put
\[
\sigma_i=a^isa^{-i},\qquad q=\sigma_{-2},\qquad
t=b^{-1}a,\qquad h_n=b^{-n}a^n.
\]
Write
\[
A(d)=\operatorname{Area}[s,\sigma_d].
\]
Conjugation, inversion, and inversion of either commutator entry preserve area. Therefore the far-commutator problem reduces to bounding \(A(d)\) for \(2\le d\le R\).

\paragraph{Six relation types for substitution}

Use the following ordered list of relation words:
\[
\begin{aligned}
F_2&=[s,\sigma_2],\\
H&=aba^{-1}b^{-1}s^{-1},\\
J_0&=bsb^{-1}\sigma_1^{-1},\\
N_2&=[\sigma_{-2},b],\\
N_3&=[\sigma_{-3},b],\\
J_1&=b\sigma_1b^{-1}\sigma_2^{-1}.
\end{aligned}
\]
Here \(J_0,J_1\) are labels for Houghton relations, unrelated to the marked central generator \(J\).

Their filling costs in the original presentation are at most
\[
(1,1,1,9,12,4).
\]

The three auxiliary costs follow from the following fillings.

For \(N_2\), two \(H\)-cells give
\[
a^2ba^{-2}=\sigma_1sb.
\]
Replacing both occurrences in \([s,a^2ba^{-2}]\) costs four \(H\)-cells. One \(J_0\)-cell turns the resulting commutator into
\[
s\sigma_1s\sigma_1^{-1}s^{-1}\sigma_1^{-1}.
\]
Three square cells turn this into \((s\sigma_1)^3\), and one braid cell fills it. Conjugation by \(a^{-2}\) proves \(N_2\) at cost \(9\).

For \(N_3\), conjugate \(N_2\) by \(a^{-1}\). In both occurrences substitute
\[
a^{-1}ba=\sigma_{-1}^{-1}b
\]
at cost two \(H\)-cells. The additional commutator
\([\sigma_{-3},\sigma_{-1}]\) costs one \(F_2\)-cell. Thus the cost is \(9+2+1=12\).

For \(J_1\), use \(ba=s^{-1}ab\) and its inverse, then \(J_0\):
\[
b\sigma_1b^{-1}
=s^{-1}a(bsb^{-1})a^{-1}s
=s^{-1}\sigma_2s.
\]
One \(F_2\)-cell finishes the transfer. The total cost is \(2+1+1=4\).

During recursion these six types remain atomic. After recursion, every auxiliary cell is expanded into its fixed original filling. The resulting count is van Kampen area in the original presentation.

\paragraph{Elementary transfer fillings}

Vectors below count the six retained cell types in the displayed order. Let \(\mathbf f_d,\mathbf n_k,\mathbf j_i,\mathbf c_m\) be filling vectors for
\[
[s,\sigma_d],\quad [\sigma_{-k},b],\quad
b\sigma_i b^{-1}\sigma_{i+1}^{-1},\quad
\sigma_m(b^msb^{-m})^{-1},
\]
respectively.

Begin with the corresponding unit vectors:
\[
\mathbf f_2=\mathbf e_{F_2},\quad
\mathbf n_2=\mathbf e_{N_2},\quad
\mathbf n_3=\mathbf e_{N_3},\quad
\mathbf j_0=\mathbf e_{J_0},\quad
\mathbf j_1=\mathbf e_{J_1}.
\]

The following constructions give componentwise upper bounds:
\begin{alignat*}{2}
\mathbf j_i
&=\mathbf j_{i-1}+2\mathbf e_H+\mathbf f_{i+1},
&&i\ge2, \tag{G.F1}\\
\mathbf c_0&=0,\qquad
\mathbf c_m=\mathbf c_{m-1}+\mathbf j_{m-1}, \tag{G.F2}\\
\mathbf f_{k+j}
&\preceq \mathbf f_k+2j\mathbf n_k+2\mathbf c_j,
&&k\ge2,\ j\ge1, \tag{G.F3}\\
\mathbf n_m
&=\mathbf n_{m-1}+\mathbf f_{m-1}+2\mathbf e_H,
&&m\ge4. \tag{G.F4}
\end{alignat*}

The recurrences follow from these substitutions:

\begin{itemize}[leftmargin=*]
\item In (G.F1), substitute \(ba=s^{-1}ab\) and its inverse, apply the \(a\)-conjugate of the preceding transfer, then commute \(s\) past \(\sigma_{i+1}\).
\item Equation (G.F2) conjugates the preceding transfer by \(b\) and applies \(J_{m-1}\).
\item For (G.F3), conjugate the target to \([\sigma_{-k},\sigma_j]\), replace \(\sigma_j\) and its inverse by \(b^jsb^{-j}\), and commute \(\sigma_{-k}\) through the \(2j\) letters \(b^{\pm1}\). The remaining commutator is a conjugate of \(F_k\).
\item For (G.F4), conjugate \(N_{m-1}\) by \(a^{-1}\), use \(a^{-1}ba=\sigma_{-1}^{-1}b\) twice, and commute \(\sigma_{-m}\) with \(\sigma_{-1}\).

\end{itemize}
Choose
\[
\begin{aligned}
\mathbf f_3&=\mathbf f_2+2\mathbf n_2+2\mathbf c_1,\\
\mathbf f_4&=\mathbf f_3+2\mathbf n_3+2\mathbf c_1,\\
\mathbf f_5&=\mathbf f_3+4\mathbf n_3+2\mathbf c_2,\\
\mathbf f_6&=\mathbf f_3+6\mathbf n_3+2\mathbf c_3.
\end{aligned}
\]
These choices are noncircular: \(\mathbf c_3\) needs \(\mathbf j_2\), which needs only \(\mathbf f_3\).

They give
\[
\begin{array}{c|rrrrrr}
&F_2&H&J_0&N_2&N_3&J_1\\ \hline
\mathbf f_2&1&0&0&0&0&0\\
\mathbf f_3&1&0&2&2&0&0\\
\mathbf f_4&1&0&4&2&2&0\\
\mathbf f_5&1&0&4&2&4&2\\
\mathbf f_6&3&4&8&6&6&4\\
\mathbf n_4&1&2&2&2&1&0\\
\mathbf n_5&2&4&6&4&3&0\\
\mathbf n_6&3&6&10&6&7&2
\end{array}
\]

\paragraph{The six-type doubling matrix}

Define a homomorphism of the free group by
\[
D(a)=a^2,\qquad D(b)=b^2,\qquad
D(s)=[a^2,b^2]=T.
\]

The free identity
\[
[a^2,b^2]
=(a[a,b]a^{-1})(ab[a,b]b^{-1}a^{-1})
[a,b](b[a,b]b^{-1})
\]
gives, with four \(H\)-cells and two \(J_0\)-cells,
\[
T=\Pi:=\sigma_1\sigma_2\sigma_0\sigma_1.
\]
Write
\[
\mathbf v=4\mathbf e_H+2\mathbf e_{J_0}.
\]

Keep this raw order, including the final positive \(\sigma_1\). The six image fillings have vectors
\[
\begin{aligned}
D(F_2):\;&4\mathbf v+\mathbf f_2+4\mathbf f_3
 +6\mathbf f_4+4\mathbf f_5+\mathbf f_6,\\
D(H):\;&0,\\
D(J_0):\;&2\mathbf v+\mathbf j_0+3\mathbf j_1
 +3\mathbf j_2+\mathbf j_3,\\
D(N_2):\;&2\mathbf v+2(2\mathbf n_3+\mathbf n_2+\mathbf n_4),\\
D(N_3):\;&2\mathbf v+2(2\mathbf n_5+\mathbf n_4+\mathbf n_6),\\
D(J_1):\;&2\mathbf v+\mathbf j_2+3\mathbf j_3
 +3\mathbf j_4+\mathbf j_5.
\end{aligned}\tag{G.F5}
\]

Here is the counting behind every row.

For \(D(F_2)=[T,a^4Ta^{-4}]\), four conversions replace \(T\) by \(\Pi\). Commuting its four factors past the four shifted factors requires sixteen far commutators. Their distances \(2,3,4,5,6\) occur with multiplicities \(1,4,6,4,1\).

The word \(D(H)\) is freely trivial by the definition of \(T\).

For \(D(J_0)\), compare \(b^2\Pi b^{-2}\) with \(a^2\Pi a^{-2}\). A factor \(\sigma_i\) uses \(J_i\) and \(J_{i+1}\). The indices \(0,1,2\) occur in \(\Pi\) with multiplicities \(1,2,1\), giving the stated sum.

For \(D(N_2)\), the shifted factors have indices
\[
-3,-2,-4,-3.
\]
Commuting each through \(b^2\) uses twice its \(N\)-filling. For \(D(N_3)\), the indices are
\[
-5,-4,-6,-5.
\]
Finally, \(D(J_1)\) uses two-step transfers on the indices \(3,4,2,3\).

Substitution of (G.F1)--(G.F4) into (G.F5) yields the matrix
\[
M=
\begin{pmatrix}
18&20&64&34&34&12\\
0&0&0&0&0&0\\
5&18&17&10&2&7\\
2&12&8&6&6&0\\
16&40&52&32&28&4\\
22&52&72&44&36&20
\end{pmatrix}. \tag{G.F6}
\]

Square and braid cells appear only when the retained auxiliary relations are expanded at the end. They do not enter the recursive substitution. Preservation of these six relation types suffices; the argument does not require a claim that \(D\) preserves every original relator.

\paragraph{An all-depth rewriting of \(h_n\)}

Put
\[
u_i=t^{-i}qt^i.
\]
The identities
\[
\sigma_{-3}=u_1,\qquad \sigma_{-4}=u_2
\]
cost, respectively,
\[
\mathbf e_{N_2},\qquad
3\mathbf e_{N_2}+\mathbf e_{N_3}.
\]

Indeed,
\[
t^{-1}qt=a^{-1}bqb^{-1}a
\]
uses \(N_2\) once. For the second identity, conjugate the first by \(a^{-1}\). Since \(a=bt\),
\[
a^{-1}u_1a=t^{-1}b^{-1}u_1bt.
\]
Commuting \(u_1\) with \(b\) costs two replacements by \(\sigma_{-3}\) and one \(N_3\)-cell. Including the initial replacement gives the stated cost.

Define formal substitutions in \(t,q\):
\[
\begin{aligned}
d(t)&=t^2q,&d(q)&=u_1qu_2u_1,\\
E(t)&=t^2qt^{-1},&E(q)&=q,\\
o&=E\circ d.
\end{aligned}
\]
They represent \(D(t),D(q)\), and their conjugates by \(b^{-1}\).

The costs and freely reduced formal lengths are
\[
\begin{array}{c|rrrrrr|r}
&F_2&H&J_0&N_2&N_3&J_1&\text{length}\\ \hline
d(t)&0&1&0&0&0&0&3\\
d(q)&0&4&2&5&1&0&10\\
o(t)&0&3&0&1&0&0&7\\
o(q)&0&10&2&9&1&0&26
\end{array} \tag{G.F7}
\]

To verify these costs, \(h_2=t^2q\) uses one \(H\)-cell. The conversion for \(D(q)\) gives
\[
\sigma_{-3}q\sigma_{-4}\sigma_{-3},
\]
so its cost is \(\mathbf v+5\mathbf e_{N_2}+\mathbf e_{N_3}\). Conjugation by \(b^{-1}\) acts on \(t\) at one \(H\)-cell and on \(q\) at one \(N_2\)-cell. The reduced word \(d(q)\) contains six signed \(t\)-letters and four signed \(q\)-letters, which gives the last row. The third row follows from the two \(t\)-letters and one \(q\)-letter of \(d(t)\).

Now set
\[
W_1=t,\qquad
W_{2m}=d(W_m),\qquad
W_{2m+1}=o(W_m)t,
\]
freely reducing after substitution.

The free identities
\[
h_{2m}=D(h_m),\qquad
h_{2m+1}=b^{-1}D(h_m)bt
\]
prove by induction that \(h_n=W_n\) has a certificate using the six retained relation types for every \(n\).

Let \(\mathbf z_n\) be its row vector of cell counts and let \(\ell_n=|W_n|\). Equations (G.F6)--(G.F7) imply
\[
\mathbf z_n\preceq
\mathbf z_{\lfloor n/2\rfloor}M
+\ell_{\lfloor n/2\rfloor}\mathbf b,
\qquad
\ell_n\le26\ell_{\lfloor n/2\rfloor}+1,
\tag{G.F8}
\]
where one may take
\[
\mathbf b=(0,10,2,9,1,0).
\]
The \(t\)-exponent of \(W_n\) is \(n\), so \(\ell_n\ge1\). Hence
\[
\ell_n\le27\ell_{\lfloor n/2\rfloor}.
\]

\paragraph{Solve the recurrence}

The characteristic polynomial is
\[
\det(tI-M)
=t\bigl(t^5-89t^4+824t^3-3416t^2-6672t+19648\bigr).
\]

Deleting the zero row and its corresponding column leaves a nonnegative \(5\times5\) matrix whose square is strictly positive. Its Perron eigenvalue \(\lambda\) is therefore simple and strictly dominates the moduli of its other eigenvalues.

An entirely integer upper certificate is
\[
w=(200,1,53,24,151,248)^T,
\]
for which
\[
Mw=(15938,0,4197,1886,11984,19720)^T\le80w.
\]
Thus \(\lambda\le80\). If \(Q\) denotes the degree-five polynomial above, then
\[
Q(79)=-5063370,\qquad Q(80)=30871488.
\]
There is an eigenvalue in \((79,80)\), so
\[
79<\lambda<80,\qquad \|M^k\|=O(\lambda^k).
\]

For \(k=\lfloor\log_2n\rfloor\), unrolling (G.F8) bounds the cell counts by a constant times
\[
\sum_{j=0}^{k-1}27^j\lambda^{k-1-j}=O(\lambda^k),
\]
because \(\lambda>27\). Expansion of the auxiliary cells costs only the fixed linear functional with coefficients \((1,1,1,9,12,4)\). Therefore
\[
\operatorname{Area}(h_nW_n^{-1})=O(n^p),
\qquad p=\log_2\lambda.
\]

Both \([s,t]\) and \([s,q]\) cost one original cell: respectively a conjugate of \(J_0^{\pm1}\) and a conjugate of \(F_2\). Replacing the two occurrences of \(h_n\) by \(W_n\), then commuting \(s\) past its letters, gives
\[
U(n):=\operatorname{Area}[h_n,s]=O(n^p). \tag{G.F9}
\]

The transfer relation
\[
\sigma_n=b^nsb^{-n}
\]
has a filling of cost \(U(n)\), since its relation word is conjugate to \([h_n,s]\). Applying (G.F3) with \(k=2,j=n\), and expanding \(N_2\) at cost \(9\), gives
\[
A(n+2)\le1+18n+2U(n).
\]
Together with (G.F9), this proves
\[
\boxed{A^*(R)=O(R^{\log_2\lambda})}.
\]

Numerically,
\[
p=6.30646825401196\ldots.
\]

\end{proof}
